\chapter[Methodology]{Methodology}
\label{cp:Methodology}

{
	\parindent0pt
	Ayang is built on a scalable architecture designed for an engaging user experience, comprising an AI model, a web application, and backend infrastructure.

	\begin{description}
		\item[Software Architecture:]

		      The user interface is a dynamic web application built with Next.js, enabling server-side rendering for fast performance even on low-bandwidth connections. The interactive map for defining community areas uses the  \href{https://www.mapbox.com}{Mapbox} API. The backend is also handled by Next.js, using Server Actions to manage user authentication, image processing, and calls to the Gemini API, which simplifies the architecture. Data is stored in Neon, a serverless Postgres database platform.


		\item[AI \& Image Analysis Pipeline:]  The AI-driven process is central to Ayang.

		      \begin{itemize}
			      \item A user uploads a "before" photo of a littered area with its GPS data.
			      \item The backend sends the image to the Gemini Vision API to identify waste types and estimate their total weight in kilograms.
			      \item After cleaning, the user uploads an "after" photo. The model performs a comparative analysis to verify the cleanup.
			      \item Upon successful verification, points are calculated based on weight and awarded to the user.
		      \end{itemize}

		\item[Model Fine-Tuning:] While the base Gemini model is capable, its accuracy can be significantly improved for niche tasks like weight estimation through fine-tuning. Our plan involves creating a custom dataset of local trash images. We will collect images of common local waste items, annotating each with a bounding box, trash type, and actual weight. This dataset will be used to fine-tune the model via Google's Vertex AI, making it more adept at recognizing local waste and estimating its weight accurately.

		\item[Community Features: ] To foster engagement, users can create communities by defining a location and radius on a map. Each community has a leaderboard to drive friendly competition. Future iterations will allow community leaders to schedule cleanup events.

	\end{description}

}